% Point to the main file
\documentclass[../../Main.tex]{subfiles}

% ========== Start the document ==========
\begin{document}

\section{Section 5.1 - Proof By Induction}

This section really only covers one thing and that proving by induction

Let's start by proving that

$$\sum_{k = 1}^{n}(2k - 1) = n^2$$

Proof 1: Break everything up

\begin{align*}
    \sum_{k = 1}^{n}(2k - 1) &= \sum_{k = 1}^{n} (2k) - \sum_{k = 1}^{n} \\
    &= 2 \sum_{k = 1}^{n} k - \sum_{k = 1}^{n} 1 \\
    &= 2 \cdot \frac{n(n + 1)}{2} - n \\
    &= n^2 + n - n \\
    \Aboxed{&= n^2} \\
\end{align*}

Proof 2: Represent as a picture (dots)

\begin{align*}
    \sum_{k = 1}^{1} (2k - 1) &= 1 = 1^2 \\
    \begin{matrix}
        \bullet
    \end{matrix} \\
    % 
    \sum_{k = 1}^{2} (2k - 1) &= 1 + 3 = 2^2 \\
    \begin{matrix}
        \bullet & \bullet \\
        \bullet & \bullet \\
    \end{matrix} \\
    % 
    \sum_{k = 1}^{5} (2k - 1) &= 1 + 3 + 5= 1^2 \\
    \begin{matrix}
        \bullet & \bullet & \bullet \\
        \bullet & \bullet & \bullet \\
        \bullet & \bullet & \bullet \\
    \end{matrix}
\end{align*}

If $\sum_{k = 1}^{n} (2k - 1) = n^2$

\[
\begin{NiceMatrix}
    \bullet & \cdots & \bullet \\
    \vdots & \ddots & \vdots \\
    \bullet & \cdots & \bullet \\
\CodeAfter
  \OverBrace{1-1}{1-3}{n}
\end{NiceMatrix}
\]

This is a square matrix so $n \times n$

Then $$\sum_{k = 1}^{n + 1} (2k - 1) = \left[ 1 + 3 + 5 + \ldots + (2n - 1)\right] + (2n + 1)$$

\[
\begin{NiceMatrix}
    \bullet & \cdots & \bullet & \bullet \\
    \vdots & \ddots & \vdots & \vdots \\
    \bullet & \cdots & \bullet & \bullet \\
    \bullet & \ldots & \bullet & \bullet \\
\CodeAfter
  \OverBrace{1-1}{1-3}{n}
\end{NiceMatrix}
\]

The extra line is $n \times 1$ and $1 \times n$, so we have $2n$ more dots, and with the extra corner dot, we get

\begin{align*}
    n^2 + 2n + 1 &= (n + 1)^2
\end{align*}

But we're just assuming that the general formula $= n^2$ is true

\begin{definition}[Proof by Induction]

\textbf{Proof by Induction} is basically proving that if a statement about the case $n$ is true, then proving the case $n + 1$ will also prove every case

Let $P(n)$ be the proposition functon 

$$P(n): \sum_{k = 1}^{n} (2k - 1) = n^2 \text{ for } n \in \setZ^+$$

\begin{center}
    Base Case: P$(1)$ is true \\
    Inductive Case: P$(n) \implies$ P$(n + 1)$
\end{center}

The above implies

\begin{align*}
    P(1) &\implies P(2) \\
    P(2) &\implies P(3) \\
    &\vdots \\
    P(n) &\implies P(n + 1) \\
\end{align*}

So we conclude that P$(n)$ is true $\forall n \in \setZ^+$

This is specifically called \textbf{Weak Induction}

\textbf{Strong Induction} is

\begin{center}
    Base Case: P$(1)$ is true \\
    Inductive Case: $P(1) \land P(2) \land \ldots \land P(n) \implies P(n + 1)$
\end{center}


\end{definition}

Let's do some examples

\begin{example}
    Prove the following

    \begin{enumerate}
        \item
        {
            $$\sum_{k = 1}^{n}k^3 = \left[ \frac{n (n + 1)}{2}\right]^2$$

            \begin{poof}
                BASE CASE:

                \begin{align*}
                    \sum_{k = 1}^{1}k^3 &= 1^3 = 1 \\
                    &= \left[ \frac{n (n + 1)}{2}\right]^2 = \left[ \frac{1(1 + 1)}{2}\right]^2 = 1 \\
                \end{align*}

                So the case $n = 1$ is true

                INDUCITVE CASE:

                Assume that $\sum_{k = 1}^{n}k^3 = \left[ \frac{n (n + 1)}{2}\right]^2$

                We then prove $\sum_{k = 1}^{n + 1}k^3 = \left[ \frac{n + 1 (n + 2)}{2}\right]^2$

                Let's do this by making the more complicated part (the proving part) equal to the simpler one

                \begin{align*}
                    \sum_{k = 1}^{n + 1}k^3 = \sum_{k = 1}^{n}k^3 + (n + 1)^3 &= 
                    \sum_{k = 1}^{n + 1}k^3 = \left[ \frac{n(n + 2)}{2}\right]^2 + (n + 1)^3 \\
                    &= \frac{n^2(n + 1)^2}{4} + (n + 1)^3 \\
                    &= (n + 1)^2 \cdot \left[ \frac{n^2}{4} + (n + 1)\right] \\
                    &= (n + 1)^2 \cdot \left[ \frac{n^2 + 4n + 4}{4} \right] \\
                    &= (n + 1)^2 \cdot \frac{(n + 2)^2}{4} \\
                    &= \frac{(n + 1)^2(n + 2)^2}{4} \\
                    \Aboxed{&= \left[\frac{(n + 1)(n + 2)}{2}\right]^2}
                \end{align*}
            \end{poof}
        }

        \item
        {
            $$\sum_{k = 1}^{n} k^2 = \frac{n (n + 1)(2n + 1)}{6} \text{ for } n \in \setZ^+$$

            \begin{poof}
                BASE CASE:

                \begin{align*}
                    \sum_{k = 1}^{1} k^2 &= 1^2 = 1 \\
                    \frac{n (n + 1)(2n + 1)}{6} &= \frac{1 \cdot (1 + 1)(2 + 1)}{6} = \frac{2 \cdot 3}{6} = 1 \\
                \end{align*}

                So out base case $n = 1$ is true
                
                INDUCTIVE CASE:

                Assume: $\sum_{k = 1}^{n} k^2 = \frac{n (n + 1)(2n + 1)}{6}$

                We prove: $\sum_{k = 1}^{n + 1} k^2 = \frac{(n + 1)(n + 2)(2n + 3)}{6}$

                Again, lets start with the more complicated looking one and make it look like the simpler one

                \begin{align*}
                    \sum_{k = 1}^{n + 1} k^2 &= \sum_{k = 1}^{n} k^2 + (n + 1)^2 \\
                    &= \frac{n (n + 1)(2n + 1)}{6} + \frac{6(n + 1)^2}{6} \\ 
                    &= (n + 1)\left[ \frac{n(2n + 1)}{6} + \frac{6(n + 1)}{6} \right] \\
                    &= (n + 1)\left[ \frac{2n^2 + n + 6n + 1}{6} \right] \\
                    &= (n + 1)\left[ \frac{2n^2 + 7n + 1}{6} \right] \\
                    &= (n + 1)\left[ \frac{(n + 2)(2n + 3)}{6} \right] \\
                    \Aboxed{&= \frac{(n + 1)(n + 2)(2n + 3)}{6}} \text{, as desired} \\
                \end{align*}
            \end{poof}
        }

        \item
        {
            $$2^n \leq n! \text{ for } n \geq 4$$

            \begin{poof}
                BASE CASE: 
                \begin{align*}
                    2^4 &= 16 \\
                    4! &= 24 \\
                    16 < 24 \\
                    \therefore 2^4 \leq 4!
                \end{align*}

                So our base case $n = 4$ is true

                INDUCTIVE CASE:

                Assume: $2^n \leq n!$

                We Prove: $2^{n + 1} \leq (n + 1)!$

                \begin{align*}
                    2^{n + 1} &= 2 \cdot 2^n \\
                    2^{n + 1} &\leq 2 \cdot n! \leq (n + 1) \cdot n! \tag{since $2 \leq (n + 1)$} \\
                    \Aboxed{&= 2^{n + 1} \leq (n + 1)!}
                \end{align*}

                \begin{note}
                    This one's kinda lame since we're dealing with inequalities
                \end{note}

            \end{poof}
        }
    \end{enumerate}
\end{example}

Let's do another proof but more words

Let's prove:

\begin{center}
    Every positive interger greater than 1 is a product of prime powers
\end{center}

\begin{poof}
    BASE CASE: $n = 2$ is a prime power
    INDUCTIVE CASE:

    Assume: $n$ is a product of prime powers

    We prove: $n + 1$ is also a product of prime powers

    Case 1:

    $n + 1$ is a prime

    Then $n + 1$ is already a prime power

    Case 2: 

    $n + 1$ is composite

    Then $n + 1$ has two factors between $1$ and $n$

    $1$ and $n$ are in our previous list which have prime powers

    $$\therefore n + 1 \text{ also does}$$
\end{poof}

\begin{example}
    Let's do an example with the Fibbonacci sequence. Recall that

    $$f_n: 1, 1, 2, \yellowbox{3}, 5, 8, 13, \yellowbox{21}, 34, 55, 89, \yellowbox{144}, 233$$

    Let's prove that every 4th Fibbonacci number is divisible by 3

    In other words,

    $$3 \mid f_{4n}, \; n \in \setZ^+$$

    \begin{poof}
        BASE CASE: $n = 1, f_4 = 3, 3 \mid 3$ so the case $n = 1$ is true

        INDUCTIVE CASE:

        Assume $3 \mid f_{4n}$

        We prove $3 \mid f_{4(n + 1)}$

        \begin{align*}
            f_{4(n + 1)} = f_{4n + 4} &= f_{4n + 3} + f_{4n + 2} \\
            &= f_{4n + 2} + f_{4n + 1} + f_{4n + 2} \\
            &= 2f_{4n + 2} + f_{4n + 1} \\
            &= 2\left( f_{4n + 1} + f_{4n} \right) + f_{4n + 1} \\
            &= 3 f_{4n + 1} + 2f_{4n} \\
        \end{align*}

        Since $3 \mid 3 f_{4n + 1}$ and $3 \mid 2f_{4n}$, then $3 \mid 3f{4n + 1} + 2f_{4n}$

        $$\therefore 3 \mid f_{4(n + 1)} \text{, as desired}$$
    \end{poof}
\end{example}

% ========== End the document ==========
\end{document}